
\chapter{Passeggiate Aleatorie (Random Walks)}

In questo capitolo esploriamo il comportamento a lungo termine di una particolare classe di processi: le \emph{passeggiate aleatorie}. 
Sia $\xi = (\xi_n)_{n \in \N}$ una successione di variabili aleatorie discrete a valori reali, e definiamo:
\[ X_0 = \xi_0, \quad X_n = X_{n-1} + \xi_n = \sum_{k=0}^n \xi_k \]
$X$ è un processo stocastico che rappresenta la posizione di una particella dopo $n$ passi di ampiezza $\xi_k$.

\section{Legge 0-1 di Kolmogorov e Lemmi di Borel-Cantelli}

Per studiare il comportamento asintotico ($n \to \infty$), è utile introdurre l'informazione portata dalla "coda" della successione.

\begin{definition}[$\sigma$-algebra coda]
Sia $\xi$ una successione di variabili. La \emph{$\sigma$-algebra coda} (tail $\sigma$-algebra) $\mathcal{T}^\xi$ è definita come l'intersezione di tutte le $\sigma$-algebre future:
\[ \mathcal{T}^\xi = \bigcap_{n \in \N} \sigma(\xi_n, \xi_{n+1}, \dots) \]
\end{definition}

Eventi che non dipendono da un numero finito di esiti iniziali (es. "la successione converge", "la somma diverge", "visita lo zero infinite volte") appartengono alla $\sigma$-algebra coda.

\begin{theorem}[Legge 0-1 di Kolmogorov]
Se le variabili $\xi_n$ sono \textbf{indipendenti}, allora ogni evento $A \in \mathcal{T}^\xi$ ha probabilità estrema:
\[ \PP(A) = 0 \quad \text{oppure} \quad \PP(A) = 1 \]
Di conseguenza, ogni variabile aleatoria misurabile rispetto a $\mathcal{T}^\xi$ è quasi certamente costante.
\end{theorem}

Per stabilire quale sia questa probabilità (0 o 1) per eventi del tipo "un evento si verifica infinite volte" (limsup di eventi), usiamo i lemmi di Borel-Cantelli.
Siano $(A_n)_{n \in \N}$ eventi e definiamo l'evento $\{A_n \text{ i.o.}\} = \limsup_{n \to \infty} A_n = \bigcap_{n} \bigcup_{k \ge n} A_k$ (i.o. sta per "infinitely often").

\begin{theorem}[Lemmi di Borel-Cantelli]
\textbf{Primo Lemma (Generale):} Se la serie delle probabilità converge:
\[ \sum_{n=1}^\infty \PP(A_n) < \infty \implies \PP(A_n \text{ i.o.}) = 0 \]
(Se un evento è "raro" in media, accadrà solo un numero finito di volte).

\textbf{Secondo Lemma (Richiede Indipendenza):} Se gli eventi $A_n$ sono indipendenti e la serie diverge:
\[ \sum_{n=1}^\infty \PP(A_n) = \infty \implies \PP(A_n \text{ i.o.}) = 1 \]
(Se tenti all'infinito un evento, anche con probabilità molto bassa, ma tale che la serie diverge, prima o poi capiterà infinite volte).
\end{theorem}

\begin{esempio}[Convergenza in probabilità ma non q.c.]
Sia $X_n$ una successione di Bernoulli indipendenti con $\PP(X_n = 1) = 1/n$.
Fissato $\epsilon \in (0, 1)$, $\PP(|X_n| > \epsilon) = 1/n \to 0$. Quindi $X_n \stackrel{P}{\to} 0$ in probabilità.
Tuttavia, $\sum \PP(X_n = 1) = \sum 1/n = \infty$. Essendo indipendenti, per il Secondo Lemma di Borel-Cantelli, $X_n = 1$ infinite volte quasi certamente. Quindi $X_n$ \emph{non converge} a 0 quasi certamente.
\end{esempio}

\section{Comportamento asintotico delle passeggiate simmetriche}

Concentriamoci su un caso specifico (modello base): i passi $\xi_n \in \{-1, 0, 1\}$ sono indipendenti e simmetrici, con $\PP(\xi_n = 1) = \PP(\xi_n = -1) = p_n \le 1/2$.

Cosa fa la particella nel lungo termine?
Usando Borel-Cantelli otteniamo la seguente dicotomia:
\begin{enumerate}
    \item Se $\sum_{n=1}^\infty p_n < \infty$, il primo lemma di Borel-Cantelli ci dice che la particella prima o poi "si ferma": da un certo $n$ in poi tutti i passi saranno nulli. (La passeggiata è definitivamente costante q.c.).
    \item Se $\sum_{n=1}^\infty p_n = \infty$, la passeggiata non si fermerà mai definitivamente. In realtà, si può dimostrare usando le martingale che in questo caso le fluttuazioni sono \textbf{illimitate}:
    \[ \limsup_{n \to \infty} X_n = +\infty \quad \text{e} \quad \liminf_{n \to \infty} X_n = -\infty \quad \text{q.c.} \]
\end{enumerate}

\section{Intermezzo: La rovina del giocatore (caso equo)}

Un'applicazione classica è la Rovina del Giocatore. 
Due giocatori (A) e (B) con capitali inziali $a, b \in \N$ scommettono 1 moneta a testa con esiti indipendenti (vince A, vince B o pareggio). Il gioco è \emph{equo}: la probabilità che vinca A in un round è uguale a quella che vinca B (chiamiamola $p > 0$).
Il gioco termina quando uno dei due finisce i soldi (capitale 0). Qual è la probabilità di rovina per A o per B?

Costruiamo una passeggiata aleatoria $X_n$ per il guadagno netto di (A). $X$ è una \emph{martingala} (poiché $\E[\xi] = p - p = 0$). Il gioco finisce al tempo d'arresto $\tau = \inf\{n : X_n = -a \text{ oppure } X_n = b\}$.
Sappiamo che la passeggiata in un gioco infinito ($\sum p = \infty$) oscilla illimitatamente, quindi $\tau < \infty$ quasi certamente (il gioco finirà sicuramente).

Applicando il teorema della martingala arrestata:
\[ \E[X_{\tau \wedge n}] = \E[X_0] = 0 \]
Usando la convergenza dominata (poiché il processo arrestato è confinato tra $-a$ e $b$), si ottiene $\E[X_\tau] = 0$.
Ma $X_\tau$ può valere solo $-a$ (se (A) si rovina, con prob. $\PP(R_A)$) oppure $b$ (se (B) si rovina, con prob. $\PP(R_B)$).
Quindi:
\[ -a \PP(R_A) + b \PP(R_B) = 0 \quad \text{e sapendo che} \quad \PP(R_A) + \PP(R_B) = 1 \]
Otteniamo la formula classica:
\[ \PP(R_A) = \frac{b}{a+b}, \quad \PP(R_B) = \frac{a}{a+b} \]

\begin{hint}[Morale]
In un gioco equo, la probabilità di sopravvivenza è strettamente proporzionale al capitale iniziale. Il giocatore con le "tasche più fonde" ha enormemente più probabilità di mandare in rovina l'avversario.
Inoltre, siccome in questo caso la passeggiata raggiunge q.c. un qualsiasi livello $\pm \alpha$, se (B) avesse un "capitale infinito" (come un Casinò), (A) andrebbe inevitabilmente in rovina con probabilità 1: $\lim_{b \to \infty} \PP(R_A) = 1$.
\end{hint}
